% Section 8: Limitations and Future Work
\section{Limitations and Future Work}
\label{sec:limitations}

Any protocol that enforces accountability architecturally---rather than by policy, convention, or organizational discipline---accepts concrete costs in exchange for that enforcement. The Recursive Spawning Protocol pays three distinct costs. Each is the direct and unavoidable price of a specific structural guarantee, and each defines a vector for continued work.

\subsection{Performance Overhead from Immutable Logging}
\label{sec:perf-overhead}

Every spawn event under the Recursive Spawning Protocol produces at minimum one Fact Layer ledger entry: a hash-chained, timestamped record that includes the child identity, the immutable parent pointer $\alpha(c) = p$, scope binding, depth value, and the Purpose Layer authorization token that justified the spawn. This logging is not optional. It is the structural mechanism that enforces chain integrity and drift prevention. The guarantees are purchased with write overhead.

At each spawn step the Fact Layer must (i) compute the hash of the preceding ledger entry, (ii) assemble the spawn record, (iii) write the record to the append-only ledger, and (iv) verify hash-chaining integrity of the resulting chain. These operations are $O(1)$ individually but compound under deep spawn chains and high-frequency spawning workloads. In the experimental evaluation (\S~\ref{sec:experimental}), chains of depth $d \geq 8$ under the immutable logging regime exhibited spawn latency overhead of $18\%$ relative to mutable-pointer baseline implementations at equivalent chain depths. The overhead grows monotonically with depth, creating a structural tension between chain length and spawn responsiveness.

The overhead is the direct price of the immutable parent pointer. We do not recommend reducing Fact Layer logging fidelity as a performance optimization. Any reduction in write completeness compromises chain integrity: a partially logged spawn chain can be retroactively manipulated in the unwritten intervals. The chain integrity theorem (\S~\ref{sec:bx3-integration}) depends on the append-only property holding for every spawn event without exception. The trade-off between accountability and write throughput is fundamental, not accidental.

Two mitigation approaches warrant future investigation. \emph{Batch hashing} groups $k$ consecutive ledger entries and commits their collective hash at intervals, reducing per-entry hash computation by a factor of $k$, at the cost of $O(k)$ ledger write latency amortized across $k$ entries. \emph{Parallel ledger writers} distribute entries across $w$ concurrent append-only stores, each with independent hash chains, merged via Byzantine fault-tolerant consensus at read time. Both approaches reduce the constant factor of the overhead; neither eliminates the fundamental trade-off.

\subsection{Scalability Under Concurrent Spawn Trees}
\label{sec:scalability}

The Recursive Spawning Protocol is specified and evaluated in the single-chain regime: one root agent producing one linear spawn chain, with each child spawning at most one subordinate. This regime is clean for formal analysis and for the proofs in \S~\ref{sec:rs-protocol}. Production multi-agent systems, however, produce concurrent, branching spawn trees where many agents at the same depth level spawn children simultaneously. Under concurrent branching, three scalability pressures emerge that the single-chain analysis does not capture.

First, the Fact Layer becomes a contention point. When $n$ concurrent child spawns target different parent nodes simultaneously, the Fact Layer pre-commit hook must serialize their write ordering to maintain hash-chain continuity. In the worst case, $n$ concurrent spawns at depth $d$ produce $O(n)$ serialized ledger writes before any children at depth $d+1$ can be activated. A spawn tree of width $w$ and depth $d$ has $w^d$ total nodes; under serialized Fact Layer writes, activation latency for level $d$ alone is $O(w^d)$, degrading rapidly as width grows.

Second, the scope refinement constraint becomes a coordination problem across sibling agents. When two concurrent child agents each hold a subset of the parent's scope, and both attempt to spawn grandchildren with scope subsets of their respective holdings, the Fact Layer pre-commit hook must evaluate each spawn independently and sequentially. A more scalable approach would require a scope allocation protocol that pre-reserves scope partitions for concurrent sibling chains, preventing redundant scope subset evaluations.

Third, the human anchor $h(n)$ becomes a sequential bottleneck. The escalation path terminates exclusively at the human anchor, which can issue only one directive at a time. Under high-concurrency spawn trees with many simultaneous exception conditions, the escalation queue at the human anchor can grow faster than the anchor can process it, limiting spawn tree width in any deployment where exceptions are frequent. Future work should evaluate hierarchical escalation trees: where $h(n)$ delegates escalation authority to trusted sub-agents who process exception conditions in parallel on its behalf, within bounded scope, and report consolidated summaries upward. This retains the exclusive human terminus property while reducing the sequential bottleneck.

\subsection{Compromised Purpose Layer}
\label{sec:purpose-compromise}

The Recursive Spawning Protocol provides strong guarantees against drift \emph{within a chain initialized from an honest, well-formed root}. It does not and cannot guarantee the honesty of the root or the inviolability of the Purpose Layer implementation. This is not a gap in the protocol specification; it is a fundamental limitation of any accountability architecture that takes initial authority as an axiomatic input.

Consider two compromise scenarios. First, a \emph{colluding root}: a human principal who deliberately initializes a root agent with a purpose clause whose scope is broader than the principal's legitimate authority. The chain propagates this over-authorized purpose faithfully; every child inherits a properly scoped subset of an improperly broad root purpose. The drift prevention property holds within the corrupted chain because each child scope is a legitimate subset of its parent's scope. But the chain operates outside any legitimate authority boundary. The Recursive Spawning model provides no defense against this scenario, because it trusts the root purpose clause as an axiomatic input.

Second, a \emph{Purpose Layer exploit}: an attacker who compromises the Purpose Layer's purpose clause registry and retroactively modifies the purpose clause of an active root agent. The modification is written to the Fact Layer as a state transition event; the chain continues operating under the new purpose clause. The immutable parent pointer chain remains intact; chain integrity is satisfied; drift prevention continues to hold because each child scope is still a subset of the now-maliciously altered parent scope. The entire accountability guarantee rests on the Purpose Layer registry being tamper-evident. If the registry is compromised, the RSP guarantees hold for a corrupted chain.

These scenarios isolate the Recursive Spawning Protocol's precise trust assumption: the model trusts the root purpose clause and the Fact Layer write protocol. Everything else follows structurally. Hardening the Purpose Layer---through hardware security modules (HSMs) that enforce purpose clause immutability at the hardware level, formal verification of the Fact Layer pre-commit hook implementation, and runtime attestation that verifies the Purpose Layer registry state against a known-good reference---is the direct and necessary continuation of this work. The trust assumption is not diffuse; it is precisely located, which makes it addressable.
